% Options for packages loaded elsewhere
\PassOptionsToPackage{unicode}{hyperref}
\PassOptionsToPackage{hyphens}{url}
\PassOptionsToPackage{dvipsnames,svgnames,x11names}{xcolor}
\documentclass[
  10pt,
  fontset=fandol]{ctexart}
\usepackage{xcolor}
\usepackage[margin=16mm]{geometry}
\usepackage{amsmath,amssymb}
\setcounter{secnumdepth}{-\maxdimen} % remove section numbering
\usepackage{iftex}
\ifPDFTeX
  \usepackage[T1]{fontenc}
  \usepackage[utf8]{inputenc}
  \usepackage{textcomp} % provide euro and other symbols
\else % if luatex or xetex
  \usepackage{unicode-math} % this also loads fontspec
  \defaultfontfeatures{Scale=MatchLowercase}
  \defaultfontfeatures[\rmfamily]{Ligatures=TeX,Scale=1}
\fi
\usepackage{lmodern}
\ifPDFTeX\else
  % xetex/luatex font selection
\fi
% Use upquote if available, for straight quotes in verbatim environments
\IfFileExists{upquote.sty}{\usepackage{upquote}}{}
\IfFileExists{microtype.sty}{% use microtype if available
  \usepackage[]{microtype}
  \UseMicrotypeSet[protrusion]{basicmath} % disable protrusion for tt fonts
}{}
\makeatletter
\@ifundefined{KOMAClassName}{% if non-KOMA class
  \IfFileExists{parskip.sty}{%
    \usepackage{parskip}
  }{% else
    \setlength{\parindent}{0pt}
    \setlength{\parskip}{6pt plus 2pt minus 1pt}}
}{% if KOMA class
  \KOMAoptions{parskip=half}}
\makeatother
\setlength{\emergencystretch}{3em} % prevent overfull lines
\providecommand{\tightlist}{%
  \setlength{\itemsep}{0pt}\setlength{\parskip}{0pt}}
\usepackage{amsmath,amssymb,mathtools,needspace}
\xeCJKDeclareCharClass{CJK}{"2032,"0400->"04FF,"2160->"216B,"2460->"2473,"25A0,"25C7,"25CB,"25CF}
\IfFontExistsTF{Arial Unicode MS}{\xeCJKsetup{AutoFallBack=true}\setCJKfallbackfamilyfont{\CJKrmdefault}{Arial Unicode MS}}{}
\setcounter{secnumdepth}{0}
\setlength{\emergencystretch}{2em}
\allowdisplaybreaks
\usepackage{bookmark}
\IfFileExists{xurl.sty}{\usepackage{xurl}}{} % add URL line breaks if available
\urlstyle{same}
\hypersetup{
  pdftitle={Newton 下山法},
  colorlinks=true,
  linkcolor={Maroon},
  filecolor={Maroon},
  citecolor={Blue},
  urlcolor={Blue},
  pdfcreator={LaTeX via pandoc}}

\title{Newton 下山法}
\author{}
\date{}

\begin{document}
\maketitle

\subsubsection{6.3.5 Newton 下山法}\label{newton-ux4e0bux5c71ux6cd5}

前面已经讨论过 Newton 法的局部收敛性。一般来说，Newton 法的收敛性依赖

\href{../../source-images/pdf-167.jpeg}{原页}

于初值 \(x_0\) 的选取，如果 \(x_0\) 偏离所求的根 \(x^*\) 比较远，则 Newton 法可能发散。

例如，用 Newton 法求方程

\[
x^3-x-1=0,
\tag{6.3.9}
\]

在 \(x=1.5\) 附近的一个根 \(x^*\)。设取迭代初值 \(x_0=1.5\)，用 Newton 公式

\[
x_{k+1}=x_k-\frac{x_k^3-x_k-1}{3x_k^2-1}
\tag{6.3.10}
\]

计算，得

\[
x_1=1.347\,83,\quad x_2=1.325\,20,\quad x_3=1.324\,72.
\]

迭代三次得到的结果 \(x_3\) 有六位有效数字。

但是，如果改用 \(x_0=0.6\) 作为迭代初值，则依 Newton 公式（6.3.10）迭代一次得

\[
x_1=17.9.
\]

这个结果反而比 \(x_0=0.6\) 更偏离了所求的根 \(x^*=1.324\,72\)。

为了防止迭代发散，对迭代过程再附加一项要求，即具有单调性

\[
|f(x_{k+1})|<|f(x_k)|.
\tag{6.3.11}
\]

满足这项要求的算法称为下山法。

将 Newton 法与下山法结合起来使用，即可在下山法保证函数值稳定下降的前提下，用 Newton 法加快收敛速度。为此，将 Newton 法的计算结果 \(\bar{x}_{k+1}=x_k-\dfrac{f(x_k)}{f'(x_k)}\) 与前一步的近似值 \(x_k\) 适当加权平均作为新的改进值，即

\[
x_{k+1}=\lambda\bar{x}_{k+1}+(1-\lambda)x_k,
\tag{6.3.12}
\]

其中 \(\lambda\ (0<\lambda\leqslant1)\) 称为下山因子。在希望挑选下山因子时，希望使单调性条件（6.3.11）成立。

下山因子的选择是个逐步探索的过程。设从 \(\lambda=1\) 开始反复将 \(\lambda\) 减半进行试算，如果能定出值 \(\lambda\) 使单调性条件（6.3.11）成立，则称``下山成功''。与此相反，如果在上述过程中找不到使条件（6.3.11）成立的下山因子 \(\lambda\)，则称``下山失败''，这时需另选初值 \(x_0\) 重算。

\end{document}
